% ===== 原文 pp.72--76 =====
% [原文 p.72]
\begin{equation}
  \omega(g_1,g_2,g_3)\in\C^\times\,, \tag{4.52}
\end{equation}
它（在已计入抵消项之后）刻画 $G^{(0)}$ 0-形式对称性的 't Hooft 反常。这些数可以汇总为一个函数
\begin{equation}
  \omega:\ (G^{(0)})^3\to\C^\times\,. \tag{4.53}
\end{equation}
这样的函数称为 $G^{(0)}$ 上取值于 $\C^\times$ 的 3-群上链。

\begin{definitionbox}{定义 4.3：群上链}
设 $M$ 是一个群。$G^{(0)}$ 上取值于 $M$ 的 $p$-群上链 $\alpha_p$ 是一个函数
\begin{equation}
  \alpha:\ (G^{(0)})^p\to M\,. \tag{4.54}
\end{equation}
\end{definitionbox}

有一个称为\textbf{五边形恒等式}（pentagon identity）的相容性条件作用在 $\omega$ 上，它来自令图 16 所示的两组移动相等
\begin{equation}
  \delta\omega(g_1,g_2,g_3,g_4)=\frac{\omega(g_2,g_3,g_4)\omega(g_1,g_2g_3,g_4)\omega(g_1,g_2,g_3)}{\omega(g_1g_2,g_3,g_4)\omega(g_1,g_2,g_3g_4)}=1\,, \tag{4.55}
\end{equation}
这意味着 $\omega$ 是 $G^{(0)}$ 上的闭 3-群上链。

\begin{definitionbox}{定义 4.4：群上链的微分与群上同调}
微分 $\delta$ 把 $\alpha_p$ 映为 $G^{(0)}$ 上取值于 $M$ 的 $(p+1)$-群上链
\begin{equation}
\begin{gathered}
  \delta\alpha_p(g_1,g_2,\cdots,g_{p+1})=\alpha_p(g_2,g_3,\cdots,g_{p+1})\alpha_p^{s(p+1)}(g_1,g_2,\cdots,g_p)\\
  \prod_{i=1}^{p}\alpha_p^{s(i)}(g_1,g_2,\cdots,g_{i-1},g_ig_{i+1},g_{i+2},g_{i+3}\cdots,g_{p+1})\,,
\end{gathered}
\tag{4.56}
\end{equation}
其中 $s(n)=(-1)^n$，即对偶数 $i$ 取 $\alpha_p$，对奇数 $i$ 取 $\alpha_p^{-1}$。

使用此微分得到的第 $p$ 阶上同调群记作
\begin{equation}
  H^p(G^{(0)},M)\,, \tag{4.57}
\end{equation}
并称为 $G^{(0)}$ 上取值于 $M$ 的第 $p$ 阶\textbf{群上同调}。
\end{definitionbox}

% [Tutor Note: p.72 式 (4.56) 是标准的非齐次群上链微分：乘积对 $i=1,\dots,p$ 求和，第 $i$ 项把 $g_i$ 与 $g_{i+1}$ 合并、其余下标依次为 $g_1,\dots,g_{i-1},g_ig_{i+1},g_{i+2},\dots,g_{p+1}$。原文的 $g_{i+2},g_{i+3}\cdots,g_{p+1}$ 正是 $g_{i+2},\dots,g_{p+1}$ 的展开写法，不是笔误。可验证：取 $p=3$、$M=\C^\times$ 展开，领头项为 $\omega(g_2,g_3,g_4)\omega(g_1,g_2,g_3)$，乘积项 $i=1,2,3$ 分别给出 $\omega(g_1g_2,g_3,g_4)^{-1}$、$\omega(g_1,g_2g_3,g_4)^{+1}$、$\omega(g_1,g_2,g_3g_4)^{-1}$，正好重现原文的 (4.55)。]

添加抵消项对应于如下改变联结处的局域算符
\begin{equation}
  \cO(g_1,g_2)\to\beta(g_1,g_2)\cO(g_1,g_2)\,,\qquad \beta(g_1,g_2)\in\C^\times\,, \tag{4.58}
\end{equation}
它把缔合子修改为
\begin{equation}
  \omega\to\omega\times\delta\beta\,, \tag{4.59}
\end{equation}

% [原文 p.73]
\begin{figure}[htbp]
\centering
\includegraphics[width=0.5\textwidth]{fig16.png}
\figtitle{16}{通过连续执行 F-移动，我们可以用两种不同方式在同一初末构型之间过渡。令这两种方式产生的 $\omega$ 因子相等，就得到式 (4.55)}
\end{figure}

% [原文 p.74]
正如读者可以核验的那样。因此，只有 $\omega$ 的上同调类 $[\omega]$ 才决定一个 't Hooft 反常。

总之，$2d$ 中离散 0-形式对称群 $G^{(0)}$ 的 't Hooft 反常由其取值于 $\C^\times$ 的第 3 阶群上同调的元素刻画
\begin{equation}
  [\omega]\in H^3(G^{(0)},\C^\times)\,. \tag{4.60}
\end{equation}
为后文之便，我们指出
\begin{equation}
  H^p(G^{(0)},\C^\times)=H^p(G^{(0)},U(1))\,, \tag{4.61}
\end{equation}
因此 $[\omega]$ 可以看作 $H^3(G^{(0)},\U(1))$ 的元素。

% [Tutor Note: p.74 式 (4.61) 的成立理由是 $\C^\times\simeq \U(1)\times\R$，且对有限群 $G^{(0)}$ 有 $H^p(G^{(0)},\R)=0$（$\R$ 可除且无挠，故上同调消失）；于是可取 $[\omega]\in H^3(G^{(0)},\U(1))$。可用 $\Z_2$ 核对：$H^1(\Z_2,\U(1))=\Z_2$、$H^2(\Z_2,\U(1))=0$、$H^3(\Z_2,\U(1))=\Z_2$，即奇阶为 $\Z_2$、偶阶为 $0$（由泛系数定理 $H^n(G,\U(1))\cong\operatorname{Hom}(H_n(G,\Z),\U(1))$，因 $\U(1)$ 可除故 $\operatorname{Ext}$ 项消失）。]

\subsubsection{SPT 相}

't Hooft 反常与 SPT 相的概念紧密相连，而 SPT 相在理论物理的许多不同领域中都非常重要。我们先给出 hep-th 文献中常用的对称性保护拓扑（symmetry protected topological，SPT）相的定义

\begin{definitionbox}{定义 4.5：SPT 相}
一个受一组 $p\ge0$-形式对称性
\begin{equation}
  \cS=\{G^{(p_1)},G^{(p_2)},\cdots,G^{(p_n)}\}\,, \tag{4.62}
\end{equation}
保护的 $d$ 维 SPT 相，是一个允许这些对称性 $\cS$ 的可逆 $d$ 维拓扑量子场论（TQFT），并且当所有背景场都被关闭时其配分函数为平凡
\begin{equation}
  Z[B_{p_1+1}=0,B_{p_2+1}=0,\cdots,B_{p_n+1}=0]=1\,. \tag{4.63}
\end{equation}
对一般的背景场，这些配分函数是相位因子
\begin{equation}
  Z[B_{p_1+1},B_{p_2+1},\cdots,B_{p_n+1}]\in U(1)\,. \tag{4.64}
\end{equation}
\end{definitionbox}

\para{忠实与非忠实对称性。} 由于没有对称性背景时配分函数全都平凡，这意味着在不考虑对称性时，SPT 相与平凡理论相同。平凡理论中唯一的算符是恒等算符
\begin{equation}
  D^{\mathrm{id}}_q\,, \tag{4.65}
\end{equation}
其维数一般地满足 $0\le q\le d-1$。因此，构成对称性 $G^{(p_i)}$ 的拓扑算符必须全都与恒等算符相同
\begin{equation}
  U_g=D^{\mathrm{id}}_{d-p_i-1}\,,\qquad \forall g\in G^{(p_i)},\ \forall i\,. \tag{4.66}
\end{equation}

% [原文 p.75]
因此，SPT 相的 $p$-形式对称性，与我们在全文中所讨论的 $p$-形式对称性相当不同：在后者中，我们一直隐含地假定没有两个拓扑算符是相同的
\begin{equation}
  U_g\neq U_{g'}\,,\qquad \text{若}\ g\neq g'\,. \tag{4.67}
\end{equation}
如果相应的拓扑算符与恒等算符相同，我们就说该对称性是\textbf{非忠实的}（non-faithful）。在这个意义上，本文一直在研究完全忠实的对称群，其中只有恒等对称性是非忠实的。另一方面，SPT 相所涉及的对称群是完全非忠实的，意思是所有对称性都是非忠实的。一般地，可以考察只有某些真子群为非忠实的群。

我们要强调，对称性的忠实性并不是一个 RG 不变的概念。从 UV 中完全忠实的一个 $p$-形式对称性 $G^{(p)}$ 出发，有可能得到一个 $G^{(p)}$ 完全非忠实的 IR 理论。下面用一个例子来说明。

\begin{examplebox}{例 4.7：从忠实到非忠实：4d $\cN=1$ $\SU(2)$ 超 Yang--Mills 理论}
考虑 4d 中规范群为
\begin{equation}
  \cG=SU(2)\,. \tag{4.68}
\end{equation}
的最小超对称纯 Yang--Mills 理论。该理论中所有物质场都在 $\SU(2)$ 的伴随表示下变换，因此该理论具有电 1-形式对称性
\begin{equation}
  G^{(1)}=Z(\cG)=\Z_2\,. \tag{4.69}
\end{equation}
该对称性非平凡地作用在 Wilson 线算符
\begin{equation}
  \cW_{j=\frac{1}{2}}\,, \tag{4.70}
\end{equation}
上，后者在 $\SU(2)$ 的自旋-$1/2$ 表示下变换。由于恒等算符必须平凡地作用在所有算符上，我们断定这个 $\Z_2$ 1-形式对称性必须是完全忠实的。

在 IR 中，该理论有两个简并真空
\begin{equation}
  v_s\,,\qquad s\in\{0,1\}\,, \tag{4.71}
\end{equation}
而每个真空中的 IR 有效理论都是 $\Z_2$ 1-形式对称性的一个 SPT 相。更确切地说，这些 SPT 的配分函数为
\begin{equation}
  Z[B_2]=\exp\left(\pi i s\int\frac{\cP(B_2)}{2}\right)\,,\qquad s\in\{0,1\}\,, \tag{4.72}
\end{equation}

% [原文 p.76]
其中 $\cP$ 是第 3.193 式附近讨论过的 Pontryagin 平方运算。因此，(4.69) 的 1-形式对称性在 IR 中是完全非忠实的。
\end{examplebox}

% [Tutor Note: 例 4.7 的方框在原文中跨 p.75--76 断开，末句 “where $\cP$ is the Pontryagin square operation...” 位于 p.76 页顶；已按规范合并为一个 examplebox。]

% [Tutor Note: p.75 式 (4.72) 中的 $\cP$ 是第 3.193 式附近的 Pontryagin 平方：它把 $\Z_2$ 值的 2-上链映为 $\Z_4$ 值的 4-上链，故 $\int\cP(B_2)$ 只在模 $4$ 意义下良定义，相位 $\pi i s\int\cP(B_2)/2$ 正是由这一模 $4$ 数据给出；$s=0,1$ 对应两个真空。（1-形式对称性的背景是 2-上链，其分类空间为 $B^2\Z_2=K(\Z_2,2)$；其 $H^4$ 取值需按系数约定复核，此处不作数值断言。）]

\para{SPT 相的信息。} 根据上述讨论，对 SPT 相而言，生成对称性的拓扑缺陷全都是平凡的。这随即引出一个问题：关于非平凡配分函数的信息究竟存储在哪里？这些信息完全存储在这些拓扑算符的联结处。

让我们用一个例子来说明这一点：$d=2$ 中离散 $G^{(0)}$ 0-形式对称群的 SPT 相。我们有拓扑线算符
\begin{equation}
  L_g=D^{\mathrm{id}}_1\,,\qquad \forall g\in G^{(0)}\,. \tag{4.73}
\end{equation}
如 (4.51) 那样，我们须在 $L_{g_1}$、$L_{g_2}$ 与 $L_{g_1g_2}$ 的联结处选择拓扑局域算符 $\cO_{g_1,g_2}$。由于这些线全都是恒等算符，联结处的局域算符是从平凡理论的真局域算符集合中选取的，而该集合是 $\C^\times$
\begin{equation}
  \cO_{g_1,g_2}=\beta(g_1,g_2)\in\C^\times\,. \tag{4.74}
\end{equation}
因此，联结处局域算符的选取被编码为 $G^{(0)}$ 上取值于 $\C^\times$ 的 2-群上链 $\beta$。要求 0-形式对称性无反常，意味着
\begin{equation}
  \delta\beta=\omega=1\,, \tag{4.75}
\end{equation}
其中 $\omega$ 已在图 15 中描述。换言之，$\beta$ 是一个上闭链。

不难看出，对给定的 $\beta$，在存在 $G^{(0)}$ 的背景 $A_1$ 时的配分函数可以表示为
\begin{equation}
  Z[A_1]=\int A_1^*\beta\,. \tag{4.76}
\end{equation}

\begin{notebox}{群上链的拉回}
给定 $G^{(0)}$ 0-形式对称性的背景场 $A_1$ 以及取值于 $M$ 的 $p$-群上链 $\alpha_p$，我们可以定义时空上一个取值于 $M$ 的 $p$-上链
\begin{equation}
  A_1^*\alpha_p\,, \tag{4.77}
\end{equation}
它称为 $\alpha_p$ 在背景 $A_1$ 下的拉回（pullback）。其定义为
\begin{equation}
  A_1^*\alpha_p(v_0,v_1,\cdots,v_p):=\alpha_p\left(A_1(v_0,v_1),A_1(v_1,v_2),\cdots,A_1(v_{p-1},v_p)\right)\,. \tag{4.78}
\end{equation}
\end{notebox}

% [Tutor Note: p.76 的 (4.76) 与 (4.78) 需配套读：$A_1$ 把单形映射到群元素，$\beta$ 作用在 $p$ 个群元素的元组上，所以 $A_1^*\beta$ 是通常的 $\U(1)$ 值 $p$-上链；$(v_0,\dots,v_p)$ 是 $p$-单形的顶点，相邻顶点的 $A_1$ 值正是 (4.78) 中论证的来源。]
